\subsection{A quartic lift of the half-Kummer bit}

\begin{statement}{D23.C36}{Proved}{The binary tetrahedral evaluation is a character modulo four}
\end{statement}
Put
\[
 S=\begin{pmatrix}0&1\\-1&0\end{pmatrix},\qquad
 R=\begin{pmatrix}0&1\\-1&1\end{pmatrix}
                 \quad\text{in }SL_2(\mathbb F_3).
\]
For every actual multiplier-one gentle element $h=(1,f_h)$,
there is a unique $\vartheta(h)\in\mathbb Z/4$ such that
\[
                 f_h(S,R)=S^{\vartheta(h)}.
\tag{D23.C36.1}
\]
It is a continuous additive character, and
\[
             \vartheta(h)\pmod2=\Psi_3(h)/2\pmod2.
\tag{D23.C36.2}
\]
In particular the parity of the four necessary $SL_2(3)$
value tables in the root computation is the already defined
half-character, rather than an independent freely removable
arithmetic phase.

For the normalized main existence problem, one can remove
the even central phase: if a mixed actual full-GT element
exists, then one exists in the finite normal form $(1,0)$
of C34 with
\[
                       \vartheta\in\{0,1\}\subset\mathbb Z/4.
\tag{D23.C36.3}
\]
This does not remove the odd possibility.

\begin{proof}
Use the natural quotient
$SL_2(3)\twoheadrightarrow PSL_2(3)=A_4$ and its faithful
projective action on $\mathbb P^1(\mathbb F_3)$. Write
$\mathbb F_4=\mathbb F_2[\omega]$ with
$\omega^2+\omega+1=0$. Identify the four projective points
with the affine line over $\mathbb F_4$ by
\[
 \infty\longmapsto\omega^2,\quad
 0\longmapsto1,\quad
 1\longmapsto\omega,\quad
 2\longmapsto0.
\tag{D23.C36.4}
\]
A direct calculation gives
\[
             R:t\longmapsto\omega t,
             \qquad S:t\longmapsto t+\omega.
\]
Here $c_0=(1-\omega)^{-1}=\omega$. The supplied A.89--A.90,
or the exact affine calculation in the proof of A.106,
therefore gives
\[
 f_h(R,S):t\longmapsto t+\omega b,
 \qquad b=\Psi_3(h)/2\pmod2.
\]
The two-cycle reverses the pair and inverts this translation;
in characteristic two inversion does not change it. Hence
$f_h(S,R)$ projects to $S^b$ in $A_4$. The full inverse image
of $\{1,\overline S\}$ is exactly the cyclic group
$\langle S\rangle$ of order four. This proves existence,
uniqueness and (D23.C36.2), without any finite census or
pentagon assumption.

For additivity use the exact GT law
$f_{hk}=f_h\varphi_h(f_k)$, where
$\varphi_h(x)=x$ and $\varphi_h(y)=f_h^{-1}yf_h$.
At the marked pair $(S,R)$ this action is
\[
 S\longmapsto S,\qquad
 R\longmapsto S^{-\vartheta(h)}RS^{\vartheta(h)}.
\]
Evaluation of $f_k$ is thus conjugation of
$S^{\vartheta(k)}$ by $S^{-\vartheta(h)}$, which leaves it
unchanged. Consequently
$f_{hk}(S,R)=S^{\vartheta(h)+\vartheta(k)}$.
Continuity follows from evaluation in a finite group.

For the last claim choose the actual arithmetic element
$\sigma\in G_{\mathbb Q^{\mathrm{ab}}}$ with
$\rho(\sigma)=0\pmod3$ and
$a_2(\sigma)=a_3(\sigma)=1$ supplied by H2. On the Galois
subgroup A.88 identifies $\Psi_3$ with the Kummer character
of $3$, up to sign. Its half modulo two is the action on
$\sqrt[4]{3}$ over $\mathbb Q^{\mathrm{ab}}$. This is the
same bit as the action on $(3/4)^{1/4}$ because
$\sqrt2\in\mathbb Q^{\mathrm{ab}}$, so A.12 and H2 give
$\Psi_3(\sigma)/2=1\pmod2$. Therefore
$\vartheta(\sigma)$ is odd and
$\vartheta(\sigma^2)=2\pmod4$, while $\sigma^2$ has all
three normalized residues $\rho\bmod3,a_2,a_3$ equal to zero.

Start with an actual normalized mixed element $h$, and
choose $e\in\{0,1\}$ such that
$\vartheta(h)+2e\in\{0,1\}$ modulo four. Then
$k=h\sigma^{2e}$ is still normalized and mixed, with
$\vartheta(k)=\vartheta(h)+2e$. By C34 and C30, the cube
of its finite paired action is the involutive form $(1,0)$.
The ninth power has that same finite form, since nine is
three modulo six; on the quartic character, nine is one
modulo four. Thus $k^9$ has the required form and retains
$\vartheta(k)\in\{0,1\}$. Its quadratic mixed bit is
unchanged because nine is odd. This proves (D23.C36.3).
\end{proof}

\paragraph{Exact markings and scope.}
The independent finite certificate
\path{agents/characters/code/binary_tetrahedral_half_character.py}
checks the projective marking (D23.C36.4), the cyclic preimage,
and the parity labels of the root's four tables when those
records are present. The supplied source locators are
\path{DELTA_ON_GT/sections/A.tex}, A.88--A.90 and the proof
of A.106 immediately before A.107. The GT multiplication
law is Lochak--Schneps (1997), printed p.~572, equation (2).
The statement does not identify the entire quartic character
with an integral Kummer character, and it does not claim that
the odd phase can be removed while preserving $a_2=0$.
Root whole-proof review and fresh exact replays passed; the result is proved in the stated scope.
